% TOP_PROBLEM: {"id": 4600018, "title": "Full-support invariant measures", "category_id": 11, "category": {"id": 11, "name": "dynamical_systems", "display_name": "Dynamical Systems", "description": "Problems about long-term behavior of deterministic systems, Hamiltonian dynamics, and periodic orbits.", "slug": "dynamical-systems", "order_index": 11, "created_at": "2026-07-31T15:26:25.671Z"}, "status": "partially_solved", "problem_number": "AMR-045-0018", "set_id": 13, "statusQualification": "The system X is explicitly the Ledrappier three-dot system defined immediately before source Question 14.2, rather than the earlier circle example."}
% FIRST_POSTED: 2026-10-01
% ENTRY_KIND: statement_only
% UnsolvedMath ID 4600018; source catalogue code AMR-045-0018
% !TeX program = lualatex
% Definition and sources only; this file is not an LLM solution attempt.
\documentclass[11pt,a4paper]{article}
\usepackage[margin=25mm]{geometry}
\usepackage{fontspec}
\setmainfont{lmroman10-regular.otf}[BoldFont=lmroman10-bold.otf,
 ItalicFont=lmroman10-italic.otf,BoldItalicFont=lmroman10-bolditalic.otf]
\usepackage{amsmath,amssymb,mathtools}
\usepackage{unicode-math}
\setmathfont{latinmodern-math.otf}
\AtBeginDocument{\let\setminus\smallsetminus}
\usepackage{xurl,hyperref}
\hypersetup{colorlinks=true,urlcolor=blue}
\setcounter{secnumdepth}{0}
\setlength{\parindent}{0pt}
\setlength{\parskip}{5pt}
\setlength{\emergencystretch}{3em}
\begin{document}
\section*{Full-support invariant measures}\label{4600018}
{\small UnsolvedMath ID 4600018 · AMR-045-0018 · Dynamical Systems · AMR Open Problem Lists}
\subsection{Definitions and mathematical statement}
Let $\mathbb F_2=\{0,1\}$ with addition modulo two, and define the
Ledrappier three-dot system
\[
 X=\{x\in\mathbb F_2^{\mathbb Z^2}:
       x_{i,j}+x_{i+1,j}+x_{i,j+1}=0
       \text{ for all }(i,j)\in\mathbb Z^2\}.
\]
It is a compact abelian group under coordinatewise addition and the
product topology. Let $\alpha^{(a,b)}x$ have coordinate
$(\alpha^{(a,b)}x)_{i,j}=x_{i+a,j+b}$, and let $m_X$ be normalized Haar
probability on the compact group $X$.

The question asks whether $m_X$ is the only Borel probability $\mu$ on
$X$ which has all three properties: $\alpha^v_*\mu=\mu$ for every
$v\in\mathbb Z^2$; the whole $\mathbb Z^2$ action is ergodic for $\mu$;
and $\operatorname{supp}\mu=X$. Ergodicity means that a measurable set
invariant modulo $\mu$ under every $\alpha^v$ has measure zero or one.
Full support means every nonempty relatively open cylinder in $X$ has
positive measure.

No ergodicity of each individual generator, and no positive entropy
assumption, is added. Haar probability is distinguished by invariance
under all additions by elements of $X$; a competing measure is required
to be invariant only under the spatial shifts. This explicit system is
the $X$ of source Question 14.2; it resolves the catalogue's imprecise
reference to the preceding commuting-endomorphism discussion.
\subsection{Short English statement}
Is Haar measure the sole ergodic, fully supported, shift-invariant probability on the binary three-dot system?
\subsection{Sources}
\begin{itemize}
\item[S1] ULAM.AI and the UnsolvedMath Contributors, supplied problems.json, record 4600018 (AMR-045-0018), AMR Open Problem Lists. Catalogue identity and source transcription; CC BY 4.0.
\url{https://huggingface.co/datasets/ulamai/UnsolvedMath}.
\item[S2] Boyle - Open Problems in Symbolic Dynamics (2008), Problem/Question/Conjecture 14.2. Original source indexed by AMR-045-0018.
\url{https://www.math.umd.edu/~mboyle/papers/openfinalsub3nov2007.pdf}.
\end{itemize}
\end{document}
